\documentclass[12pt,a4paper]{article}

% ── Packages ──────────────────────────────────────────────────────────────
\usepackage[utf8]{inputenc}
\usepackage[T1]{fontenc}
\usepackage{amsmath,amssymb,amsthm}
\usepackage{mathtools}
\usepackage{geometry}
\usepackage{hyperref}
\usepackage{cleveref}
\usepackage{booktabs}
\usepackage{array}
\usepackage{longtable}
\usepackage{graphicx}
\usepackage{xcolor}
\usepackage{tikz}
\usepackage{algorithm}
\usepackage{algpseudocode}
\usepackage{listings}
\usepackage{caption}
\usepackage{subcaption}
\usepackage[numbers,sort&compress]{natbib}

\geometry{margin=2.5cm}
\hypersetup{colorlinks=true,linkcolor=blue!60!black,citecolor=blue!60!black,urlcolor=blue!60!black}

% ── Theorem environments ──────────────────────────────────────────────────
\newtheorem{theorem}{Theorem}[section]
\newtheorem{lemma}[theorem]{Lemma}
\newtheorem{proposition}[theorem]{Proposition}
\newtheorem{corollary}[theorem]{Corollary}
\theoremstyle{definition}
\newtheorem{definition}[theorem]{Definition}
\newtheorem{remark}[theorem]{Remark}

% ── Macros ────────────────────────────────────────────────────────────────
\newcommand{\ZZ}{\mathbb{Z}}
\newcommand{\RR}{\mathbb{R}}
\newcommand{\GL}{\mathrm{GL}}
\newcommand{\CWS}{\mathrm{CWS}}
\newcommand{\NF}{\mathrm{NF}}
\newcommand{\polyDmax}{5}
\newcommand{\vertNmax}{64}

\lstset{
  basicstyle=\ttfamily\small,
  keywordstyle=\color{blue!70!black}\bfseries,
  commentstyle=\color{green!50!black},
  stringstyle=\color{red!60!black},
  numbers=left,
  numberstyle=\tiny\color{gray},
  frame=single,
  breaklines=true,
  captionpos=b
}

% ══════════════════════════════════════════════════════════════════════════
\title{%
  Complete Classification of Five-Dimensional Reflexive Polytopes\\
  from Weight Systems via Normal Form Hashing
}

\author{
  Andreas Hatziiliou%
  %\thanks{Email: \texttt{...}}
  \\
  Eldar Sultanow
  %\textit{Affiliation}
}

\date{\today}

% ══════════════════════════════════════════════════════════════════════════
\begin{document}
\maketitle

\begin{abstract}
We present a computational pipeline for the complete classification of
five-dimensional reflexive polytopes arising from combined weight systems.
Starting from a dataset of approximately 183~billion combined weight
systems, we compute the normal form of each associated polytope under
$\GL(5,\ZZ)$ equivalence using the PALP library, hash the resulting
canonical matrix via xxHash128, and deduplicate across the full dataset
using a two-phase external sort-merge algorithm.  The pipeline processes
the 3.2~TB input dataset across distributed compute nodes, producing a
single deduplicated catalogue of unique polytopes with frequency counts
and topological metadata.  Formal verification using bounded model
checking (CBMC) and abstract interpretation (Frama-C) establishes
correctness of all critical code paths, including sort ordering,
count preservation, binary layout, and normal form determinism.

% TODO: insert final count of unique polytopes
\end{abstract}

% ══════════════════════════════════════════════════════════════════════════
\section{Introduction}
\label{sec:intro}

The classification of reflexive polytopes is a foundational problem in
string theory and algebraic geometry.  A \emph{reflexive polytope}
$\Delta \subset \RR^d$ is a lattice polytope containing the origin in its
interior whose dual polytope $\Delta^*$ is also a lattice polytope.  By
Batyrev's construction~\cite{Batyrev1994}, each pair $(\Delta, \Delta^*)$
of reflexive polytopes determines a mirror pair of Calabi--Yau manifolds
as hypersurfaces in the associated toric varieties.  The classification
of reflexive polytopes in a given dimension therefore yields a
classification of a large family of Calabi--Yau manifolds and their
mirror partners.

The classification of reflexive polytopes has been completed in dimensions
up to four.  In dimension~2, there are exactly 16 reflexive
polygons~\cite{Koelman1991}.  In dimension~3, Kreuzer and Skarke
identified 4,319 reflexive polytopes~\cite{KreuzerSkarke1998}.  The
landmark computation of Kreuzer and Skarke in dimension~4 found
473,800,776 reflexive polytopes~\cite{KreuzerSkarke2000}, which has
served as the foundation for much subsequent work in string
phenomenology.  The classification in dimension~5 remains open and is the
subject of this paper.

The combinatorial explosion in higher dimensions is dramatic: while
dimension~4 required examining approximately $3 \times 10^8$ weight
systems, dimension~5 involves approximately $1.83 \times 10^{11}$
combined weight systems (CWS).  Each CWS defines a weighted projective
space in which a Calabi--Yau hypersurface can be embedded; computing the
associated reflexive polytope and its normal form requires lattice point
enumeration, vertex/facet discovery, and a canonical form computation
under the full $\GL(5,\ZZ)$ symmetry group.  The computational cost per
CWS varies by three orders of magnitude depending on the number of
lattice points (from $\sim 10\,\mu\text{s}$ for simple polytopes to
$\sim 50\,\text{ms}$ for polytopes with $> 50{,}000$ lattice points).

We describe a complete pipeline for this computation: a multi-threaded
C++ classifier that links PALP~\cite{PALP} as a library (rather than
invoking it as a subprocess), uses xxHash128 for $O(1)$ deduplication,
and employs an external sort-merge algorithm for the final global
deduplication across distributed compute nodes.  We provide formal
verification of all critical code paths using CBMC and Frama-C,
establishing mathematical guarantees that no records are lost, no counts
are corrupted, and the normal form computation is deterministic.

% ══════════════════════════════════════════════════════════════════════════
\section{Mathematical Background}
\label{sec:math}

\subsection{Reflexive polytopes and Calabi--Yau manifolds}

\begin{definition}
A \emph{lattice polytope} $\Delta \subset \RR^d$ is the convex hull of
finitely many points in $\ZZ^d$.  It is \emph{reflexive} if
$0 \in \mathrm{int}(\Delta)$ and the dual polytope
\[
  \Delta^* = \{ y \in \RR^d : \langle x, y \rangle \geq -1
               \;\text{for all}\; x \in \Delta \}
\]
is also a lattice polytope.
\end{definition}

By a theorem of Batyrev~\cite{Batyrev1994}, each $d$-dimensional reflexive
polytope $\Delta$ determines a $(d-1)$-dimensional Calabi--Yau
hypersurface in the toric variety $\mathbb{P}_\Delta$ associated to
$\Delta$.  The Hodge numbers of this Calabi--Yau manifold are
computable from the combinatorial data of $\Delta$ and $\Delta^*$
via Batyrev's formulae~\cite{BatyrevBorisov1996}.

\subsection{Combined weight systems}

A \emph{combined weight system} (CWS) in dimension~$d$ is a collection
of weight vectors $(w_0, w_1, \ldots, w_d) \in \ZZ_{>0}^{d+1}$ such
that $\gcd(w_0, \ldots, w_d) = 1$.  Each CWS defines a weighted
projective space $\mathbb{P}(w_0, \ldots, w_d)$, and the associated
Newton polytope is the convex hull of the lattice points
\[
  \Delta = \mathrm{Conv}\left\{
    m \in \ZZ^d : \sum_{i=0}^{d} w_i m_i = \deg,\;
    m_i \geq 0
  \right\},
\]
where $\deg = \sum_i w_i$ is the degree.

Multiple distinct CWS can generate the same reflexive polytope (up to
$\GL(d,\ZZ)$ equivalence).  The central computational challenge is
therefore to enumerate all CWS, compute the associated polytopes, and
identify which ones are isomorphic.

\subsection{Normal forms under $\GL(d,\ZZ)$}

Two polytopes $\Delta_1, \Delta_2 \subset \RR^d$ are
\emph{$\GL(d,\ZZ)$-equivalent} if there exists $A \in \GL(d,\ZZ)$ and
$v \in \ZZ^d$ such that $A \cdot \Delta_1 + v = \Delta_2$.  To test
equivalence, we use the \emph{normal form} algorithm implemented in
PALP~\cite{PALP}, which computes a canonical representative $\NF(\Delta)$
of the equivalence class.

\begin{proposition}[Normal form uniqueness]
For any two reflexive polytopes $\Delta_1, \Delta_2$:
\[
  \Delta_1 \sim_{\GL(d,\ZZ)} \Delta_2
  \iff
  \NF(\Delta_1) = \NF(\Delta_2).
\]
\end{proposition}

The normal form is represented as a $d \times n_v$ integer matrix, where
$n_v$ is the number of vertices.  The algorithm proceeds by:
\begin{enumerate}
  \item Sorting the vertex list into a canonical order.
  \item Constructing the vertex pairing matrix (VPM).
  \item Searching for the lexicographically minimal representation under
        all symmetries of the VPM.
\end{enumerate}

\subsection{Hash-based deduplication}

Rather than comparing the full normal form matrices (up to $5 \times 64
\times 8 = 2{,}560$ bytes each), we hash each normal form to a 128-bit
digest using xxHash128~\cite{xxHash}.  The collision probability for
$n$ distinct polytopes is bounded by the birthday paradox:
\[
  P(\text{collision}) \approx \frac{n^2}{2^{129}} \approx
  \frac{(5 \times 10^8)^2}{2^{129}} \approx 3.7 \times 10^{-22}.
\]
This is negligible for any practical purpose.

% ══════════════════════════════════════════════════════════════════════════
\section{Five-Dimensional Minimal Polytopes}
\label{sec:minimal-5d}

This section isolates the combinatorial input from
\citet{KreuzerSkarke1997} that controls which combined weight-system
patterns can occur in dimension~5.  We work with the notation of
their Lemma~1 and Corollary~1: a minimal polytope $M$ is described by an
ordered family of good simplices $S_1,\dots,S_\lambda$ such that the
partial unions span minimal polytopes $M_1,\dots,M_\lambda=M$.

\begin{definition}
Fix a Corollary~1 ordering $S_1,\dots,S_\lambda$ of a $5$-dimensional
minimal polytope.  Let $m_\mu$ be the number of vertices of $S_\mu$ that
do not occur in $S_1 \cup \cdots \cup S_{\mu-1}$.  We call
$(m_1,\dots,m_\lambda)$ the \emph{ordered new-vertex profile} of the
chosen decomposition.
\end{definition}

\begin{proposition}[Five-dimensional profile constraint]
\label{prop:5d-profiles}
Let $M$ be a $5$-dimensional minimal polytope and let
$(m_1,\dots,m_\lambda)$ be the ordered new-vertex profile of a Corollary~1
decomposition.  Then each $m_\mu \ge 2$ and
\[
  \sum_{\mu=1}^{\lambda} (m_\mu - 1) = 5.
\]
Consequently, after sorting the profile in weakly decreasing order one
obtains one of the seven possibilities
\[
  [6],\ [5,2],\ [4,3],\ [4,2,2],\ [3,3,2],\ [3,2,2,2],\ [2,2,2,2,2].
\]
The case $[6]$ is the simplex case, hence the single-weight-system case;
all other profiles correspond to genuine combined weight-system patterns.
\end{proposition}

\begin{proof}
Write $k_\mu$ for the number of vertices in $M_\mu$.  By
Corollary~1(a) of \citet{KreuzerSkarke1997},
$\dim M_\mu = k_\mu - \mu$.  Therefore the dimension jump at step $\mu$ is
\[
  \dim M_\mu - \dim M_{\mu-1}
  = (k_\mu - \mu) - (k_{\mu-1} - (\mu-1))
  = (k_\mu - k_{\mu-1}) - 1
  = m_\mu - 1.
\]
Summing these jumps from $M_1$ to $M_\lambda = M$ yields
\[
  5 = \dim M = \sum_{\mu=1}^{\lambda} (m_\mu - 1).
\]
Each $m_\mu$ is at least $2$, because every new simplex contributes at
least one new dimension and hence at least two vertices.  Subtracting $1$
from each entry therefore produces a partition of $5$ into positive
integers.  The only partitions of $5$ are
$5$, $4+1$, $3+2$, $3+1+1$, $2+2+1$, $2+1+1+1$, and $1+1+1+1+1$.
Adding $1$ back to each part gives exactly the seven profiles listed.
\end{proof}

\begin{lemma}[Lower-dimensional input data]
\label{lem:lower-dimensional-input}
Up to relabeling of vertices, the complete $3$-dimensional minimal-polytope
structures are
\[
  \{\{1,2,3,4\}\},\qquad
  \{\{1,2,3\},\{4,5\}\},\qquad
  \{\{1,2,3\},\{1,4,5\}\},\qquad
  \{\{1,2\},\{3,4\},\{5,6\}\},
\]
and the complete $4$-dimensional minimal-polytope structures are
\[
  \{\{1,2,3,4,5\}\},\qquad
  \{\{1,2,3,4\},\{1,2,5,6\}\},\qquad
  \{\{1,2,3,4\},\{1,5,6\}\},
\]
\[
  \{\{1,2,3,4\},\{5,6\}\},\qquad
  \{\{1,2,3\},\{4,5,6\}\},\qquad
  \{\{1,2,3\},\{4,5\},\{6,7\}\},
\]
\[
  \{\{1,2,3\},\{1,4,5\},\{6,7\}\},\qquad
  \{\{1,2,3\},\{1,4,5\},\{1,6,7\}\},\qquad
  \{\{1,2\},\{3,4\},\{5,6\},\{7,8\}\}.
\]
\end{lemma}

\begin{proof}
This is exactly the explicit $n=3$ and $n=4$ classification stated in
Section~4 of \citet{KreuzerSkarke1997}.  We have only rewritten the same
structures in the present brace notation for good simplices.
\end{proof}

\begin{lemma}[Dimension-$5$ extension rule]
\label{lem:5d-extension-rule}
Let $M$ be a $5$-dimensional minimal polytope.
\begin{enumerate}
  \item If Lemma~1 of \citet{KreuzerSkarke1997} selects $n'=4$, then
  $M$ is obtained from one of the $4$-dimensional base structures of
  \cref{lem:lower-dimensional-input} by adjoining a new good simplex
  \[
    S = R \cup \{a,b\},
  \]
  where $a,b$ are new vertices and $R$ is a subset of old vertices with
  $|R| \le 3$.
  \item If Lemma~1 selects $n'=3$, then $M$ is obtained from one of the
  $3$-dimensional base structures of
  \cref{lem:lower-dimensional-input} by adjoining a new good simplex
  \[
    S = R \cup \{a,b,c\},
  \]
  where $a,b,c$ are new vertices and $R$ is a subset of old vertices with
  $|R| \le 1$.
\end{enumerate}
Conversely, every candidate produced in this way satisfies the dimensional
constraint from Lemma~1.  It survives to a valid structure if and only if
it is not excluded by Lemma~2, i.e. if no good simplex has exactly one
vertex that belongs to no other simplex in the spanning collection.
\end{lemma}

\begin{proof}
Lemma~1 says that a non-simplex $5$-dimensional minimal polytope contains
an $n'$-dimensional minimal polytope $M'$ and one further good simplex $S$
such that
\[
  k-k' = 5-n'+1,
  \qquad
  \dim S \le n'.
\]
Since $n' \ge 5/2$, only $n'=4$ and $n'=3$ can occur.

If $n'=4$, then $k-k' = 2$, so the new simplex contributes exactly two new
vertices, say $a,b$.  Because $\dim S \le 4$, the simplex has at most five
vertices, hence at most three old ones.  Therefore
$S = R \cup \{a,b\}$ with $|R| \le 3$.

If $n'=3$, then $k-k' = 3$, so the new simplex contributes exactly three
new vertices, say $a,b,c$.  Because $\dim S \le 3$, the simplex has at
most four vertices, hence at most one old one.  Therefore
$S = R \cup \{a,b,c\}$ with $|R| \le 1$.

The converse statement about the dimensional constraint is immediate from
the same bounds.  The final sentence is exactly Lemma~2: a spanning set of
good simplices is forbidden precisely when one of the simplices has a
unique vertex not shared by any other simplex.
\end{proof}

\begin{theorem}[Complete $5$D structure list]
\label{thm:5d-structures}
Up to relabeling of the vertices, there are exactly $47$ overlap
structures for $5$-dimensional minimal polytopes.  They break up by total
vertex count as follows:
\begin{center}
\begin{tabular}{@{}cr@{}}
\toprule
Vertices $k$ & Number of structures \\
\midrule
6 & 1 \\
7 & 6 \\
8 & 25 \\
9 & 13 \\
10 & 2 \\
\bottomrule
\end{tabular}
\end{center}
Grouping the same $47$ structures by their sorted new-vertex profiles
produces the multiplicities
\[
  [6] : 1,\quad [5,2] : 2,\quad [4,3] : 4,\quad [4,2,2] : 7,
\]
\[
  [3,3,2] : 18,\quad [3,2,2,2] : 13,\quad [2,2,2,2,2] : 2.
\]
Appendix~\ref{app:5d-structures} lists one canonical Corollary~1 ordering
for each structure, together with its overlap matrix
$O_{ij} = |S_i \cap S_j|$.
\end{theorem}

\begin{proof}
By \cref{lem:lower-dimensional-input}, there are exactly four possible
$3$-dimensional base structures and nine possible $4$-dimensional base
structures.  By \cref{lem:5d-extension-rule}, every non-simplex
$5$-dimensional minimal polytope arises from one of these thirteen bases by
one of the following finite operations:
\begin{enumerate}
  \item choose a $4$D base and adjoin a simplex $R \cup \{a,b\}$ with
  $|R| \le 3$;
  \item choose a $3$D base and adjoin a simplex $R \cup \{a,b,c\}$ with
  $|R| \le 1$.
\end{enumerate}
The simplex case itself contributes the unique structure with profile
$[6]$.

Thus the problem has been reduced to a completely finite search.  For each
base structure, one runs through all allowed subsets $R$ of old vertices,
forms the corresponding candidate spanning collection of good simplices,
and discards the candidate exactly when Lemma~2 is violated.  The
surviving candidates are then quotient by the obvious notion of
isomorphism: permutation of the simplices together with relabeling of the
vertices.  Equivalently, one keeps one canonical representative of each
overlap hypergraph.

This search is complete because every $5$D minimal polytope must arise from
Lemma~1 and therefore from one of the two cases above.  It is correct
because the only combinatorial obstructions used by Kreuzer--Skarke at this
stage are precisely the dimensional bound from Lemma~1 and the uniqueness
obstruction from Lemma~2.  No surviving candidate is identified with a
different one unless their overlap hypergraphs are isomorphic, since the
canonicalization step quotients exactly by vertex relabeling and simplex
permutation.

Carrying out that finite search yields exactly the $47$ canonical overlap
structures listed in Appendix~\ref{app:5d-structures}.  Counting them by
total number of vertices gives
\[
  1,\ 6,\ 25,\ 13,\ 2
\]
for $k=6,7,8,9,10$ respectively, and counting them by sorted new-vertex
profile gives
\[
  [6] : 1,\quad [5,2] : 2,\quad [4,3] : 4,\quad [4,2,2] : 7,
\]
\[
  [3,3,2] : 18,\quad [3,2,2,2] : 13,\quad [2,2,2,2,2] : 2.
\]
The finite search itself, including these exact counts, is machine-checked
in \texttt{lean/MinimalPolytopes5D/Structures.lean} by the theorems
\texttt{canonicalFiveDimensionalStructures\_length},
\texttt{countsByVertex\_correct},
\texttt{countsByProfile\_correct}, and
\texttt{sortedProfile\_allowed}.
\end{proof}

\begin{remark}
The appendix records one admissible Corollary~1 ordering for each overlap
hypergraph.  Some hypergraphs admit more than one such ordering and hence
more than one ordered new-vertex profile.  The profile in the appendix is
therefore a canonical choice, not an invariant attached to the hypergraph
itself.  The arithmetic core of \cref{prop:5d-profiles} is formalized in
Lean~4 in the file \texttt{lean/MinimalPolytopes5D/Profiles.lean}, and the
finite recursive search from the explicit $3$D and $4$D base lists to the
full set of $47$ canonical $5$D structures is formalized in
\texttt{lean/MinimalPolytopes5D/Structures.lean}.
\end{remark}

% ══════════════════════════════════════════════════════════════════════════
\section{Input Data}
\label{sec:data}

% TODO: Fill in details about:
% - Source and provenance of the CWS dataset
% - Total number of CWS (183 billion across 4000 Parquet files, ~3.2 TB)
% - How the CWS were generated / enumerated
% - File format and schema (Parquet, 13 columns per row)
% - Pre-computed metadata columns (h11, h12, h13, dual_point_count, etc.)
% - Any filtering applied (reflexive-only subset)
% - Citation for the dataset
% - Distribution of CWS by vertex count, point count, etc.

% ══════════════════════════════════════════════════════════════════════════
\section{Computational Pipeline}
\label{sec:pipeline}

The pipeline operates in two phases: \emph{classification} (per-file
processing with local deduplication and checkpointing) and \emph{merge}
(global deduplication across all checkpoint shards via external
sort-merge).  We describe each phase in detail.

\subsection{Phase A: Classification}
\label{sec:classification}

Each input Parquet file ($\sim 46$M rows, $\sim 800$~MB) is processed
independently.  The six weight columns are extracted via Apache Arrow
columnar reads, and the resulting CWS rows are distributed to a pool
of worker threads using a dynamic work-stealing scheduler.

\begin{algorithm}[t]
\caption{Per-file classification}
\label{alg:classify}
\begin{algorithmic}[1]
\Require Parquet file $F$ with $n$ CWS rows; thread count $T$
\Ensure Global hash map $G$ updated; checkpoint written
\State $\mathit{rows} \gets \textsc{ReadParquet}(F)$
\State Allocate $T$ thread-local hash maps $L_0, \ldots, L_{T-1}$
\State $\mathit{next\_block} \gets 0$
\Comment{Atomic counter for work-stealing}
\ForAll{threads $t \in \{0, \ldots, T-1\}$ \textbf{in parallel}}
  \State $\mathit{ws} \gets \textsc{PalpWorkspaceAlloc}()$
  \Loop
    \State $b \gets \textsc{AtomicFetchAdd}(\mathit{next\_block}, 1)$
    \If{$b \geq \lceil n / B \rceil$} \textbf{break} \EndIf
    \For{$i = bB$ \textbf{to} $\min((b+1)B, n) - 1$}
      \State $\mathit{result} \gets \textsc{PalpComputeNF}(\mathit{ws}, \mathit{rows}[i].\mathit{weights})$
      \If{$\mathit{result.ok} = 0$} increment failed counter; \textbf{continue} \EndIf
      \State $h \gets \textsc{xxHash128}(\mathit{result.nf}, \mathit{result.dim} \times \mathit{result.nv})$
      \If{$h \in L_t$} $L_t[h].\mathit{count} \mathrel{+}= 1$
      \Else{} $L_t[h] \gets \textsc{NewEntry}(\mathit{rows}[i], \mathit{result})$
      \EndIf
    \EndFor
  \EndLoop
\EndFor
\For{$t = 0$ \textbf{to} $T-1$}
  \Comment{Merge under global lock}
  \State $\textsc{MergeMaps}(G, L_t)$
\EndFor
\State $\textsc{WriteCheckpoint}(G)$
\end{algorithmic}
\end{algorithm}

Each worker thread holds a private PALP workspace (allocated via
\texttt{calloc} for deterministic initialisation) and a thread-local
\texttt{std::unordered\_map}.  The work-stealing block size $B = 1024$
balances atomic counter contention against load-balancing granularity.

\paragraph{PALP as a library.}
We link PALP~\cite{PALP} directly as a C library rather than invoking it
as a subprocess (\texttt{poly.x -N}).  This eliminates process fork
overhead, text serialisation, and file I/O, increasing per-core
throughput from $\sim 6{,}000$ to $\sim 7{,}500$~CWS/s.  Two
thread-safety patches were required:
\begin{enumerate}
  \item Static scratch buffers in \texttt{Vertex.c:Make\_New\_CEqs} were
        changed to \texttt{\_\_thread} storage under
        \texttt{-DPALP\_THREADSAFE}.
  \item Global file pointers (\texttt{inFILE}, \texttt{outFILE}) were
        redirected to \texttt{/dev/null} at initialisation.
\end{enumerate}

\paragraph{Compilation and optimisation.}
PALP is compiled with \texttt{-O3 -march=native -flto -DPOLY\_Dmax=5}.
The \texttt{POLY\_Dmax=5} flag reduces all dimension-indexed arrays from
$6 \to 5$ elements, improving cache utilisation.  Tighter bounds tuned
for 5D polytopes further reduce memory footprint (e.g.,
\texttt{POINT\_Nmax} from 2,000,000 to 200,000).  These optimisations
yield an 18\% improvement in single-threaded throughput.

\paragraph{Critical bug: side effects in \texttt{assert()}.}
During development, we discovered that PALP contains numerous
\texttt{assert()} calls with side effects (counter increments, memory
allocations, function calls that modify state).  Under \texttt{-DNDEBUG}
(standard for release builds), these side effects are silently removed,
causing incorrect behaviour.  We introduced a \texttt{PALP\_FAST\_ASSERT}
macro that evaluates the expression for its side effects while
suppressing the abort on failure.  Frama-C abstract interpretation
confirmed that disabling these asserts causes at least one sure
uninitialized memory read (Vertex.c:598), while enabling
\texttt{PALP\_FAST\_ASSERT} produces zero sure alarms.

\subsection{Phase B: External sort-merge}
\label{sec:merge}

After all input files have been processed by one or more compute nodes,
the resulting checkpoint shards must be globally deduplicated.  Each shard
contains a subset of the unique polytopes seen by that runner, with
local frequency counts.  The same polytope may appear in multiple shards
if it was generated by CWS in different input files processed by
different runners.

The merge algorithm operates in two phases:

\paragraph{Phase 1: Sort.}
Each checkpoint shard is loaded into memory, sorted in-place by the
128-bit hash key using \texttt{std::sort} (parallel, $O(n \log n)$), and
written back.  The sort uses a lexicographic comparator on $(h_{\mathrm{hi}},
h_{\mathrm{lo}})$, which is a strict weak ordering (verified formally;
see~\cref{sec:verification}).  An atomic rename pattern
(\texttt{write .tmp} $\to$ \texttt{rename}) ensures crash safety.
Resumability is provided by \texttt{.sorted} marker files.

\paragraph{Phase 2: $N$-way streaming merge.}
All sorted shards are opened simultaneously via buffered readers
(\texttt{SortedBinaryReader}, 256K-record buffer per file, $\sim 18$~MB
RAM each).  A min-heap of size $N$ (number of shards) drives the merge:
at each step, the record with the smallest key is extracted, and all
records sharing the same key across all shards are coalesced by summing
their frequency counts.  The deduplicated stream is written directly to a
Parquet file with ZSTD compression.

Peak RAM during Phase~2 is $N \times 18\,\text{MB} \approx 1.2\,\text{GB}$
for 67~shards.  The total I/O is one sequential read of each sorted shard
plus one sequential write of the output Parquet file.

\begin{proposition}[Count preservation]
\label{prop:count}
Let $S_1, \ldots, S_N$ be the input shards and $O$ the output.  Then
\[
  \sum_{r \in O} r.\mathit{count}
  = \sum_{i=1}^{N} \sum_{r \in S_i} r.\mathit{count}.
\]
\end{proposition}
\begin{proof}
Each input record is consumed exactly once by the heap (the reader
visits every record, and each record enters and leaves the heap exactly
once).  When records are coalesced, their counts are summed.  The
emitted record carries the total count.  No records are dropped or
fabricated.  This property is verified formally by bounded model checking
(see~\cref{sec:verification}).
\end{proof}

\subsection{Distributed execution}
\label{sec:distributed}

The 4,000 input Parquet files ($\sim 3.2$~TB total) are partitioned
across multiple compute nodes using \texttt{-{}-start} and
\texttt{-{}-end} flags.  Each node processes its assigned range
independently, writing periodic checkpoints.  After all nodes complete,
the checkpoint shards are collected and merged in a single pass.

This embarrassingly parallel architecture scales linearly with the number
of nodes: each node's throughput is limited only by its CPU and memory.
The merge step is $I/O$-bound and completes in minutes.

% ══════════════════════════════════════════════════════════════════════════
\section{Performance}
\label{sec:performance}

\subsection{Per-CWS computation profile}

Profiling on a 100,000-CWS sample reveals that 78\% of execution time
is spent in \texttt{Find\_Equations} subroutines that repeatedly scan
the full lattice point list.  The per-CWS computation cost varies by
three orders of magnitude:

\begin{table}[h]
\centering
\caption{Per-CWS computation cost by lattice point count.}
\label{tab:cost}
\begin{tabular}{@{}lrrr@{}}
\toprule
Point count & Approx.\ time/CWS & \% of dataset & Cumulative \% \\
\midrule
$< 100$     & $\sim 10\,\mu$s   & 59\%   & 59\% \\
$100$--$500$  & $\sim 50\,\mu$s   & 34\%   & 93\% \\
$500$--$5{,}000$ & $\sim 200\,\mu$s & 6.3\%  & 99.3\% \\
$5{,}000$--$50{,}000$ & $\sim 5\,$ms & 0.6\%  & 99.9\% \\
$> 50{,}000$ & $\sim 50\,$ms     & 0.03\% & $\sim 100$\% \\
\bottomrule
\end{tabular}
\end{table}

This heavy-tailed distribution motivates the dynamic work-stealing
scheduler: static partitioning would cause threads assigned to expensive
polytopes to bottleneck the entire computation.

\subsection{Throughput and scaling}

On a single AMD Ryzen 9 7950X3D (16 cores / 32 threads, 96~GB RAM),
the classifier achieves $\sim 241{,}000$~CWS/s using all 32 logical
threads, representing a $\sim 14\times$ speedup over single-threaded
execution.  Sub-linear scaling is attributable to memory bandwidth
contention and L3 cache sharing across the two CCDs.

\begin{table}[h]
\centering
\caption{Classification throughput on AMD Ryzen 9 7950X3D.}
\label{tab:throughput}
\begin{tabular}{@{}lrrr@{}}
\toprule
Configuration & Throughput (CWS/s) & Wall time (46M CWS) & Speedup \\
\midrule
1 thread                & 1,730  & 26,780\,s & $1\times$ \\
32 threads, work-stealing & 241,000 & 192\,s  & $14\times$ \\
\bottomrule
\end{tabular}
\end{table}

\subsection{Deduplication statistics}

Single-file analysis reveals significant structure in the duplication
pattern:

\begin{table}[h]
\centering
\caption{Deduplication rates by vertex count (single file, 46M CWS).}
\label{tab:dedup}
\begin{tabular}{@{}rrrr@{}}
\toprule
Vertex count & Input CWS & Unique polytopes & Dedup rate \\
\midrule
6  & 327,520   & 2,862     & 99.1\% \\
7  & 1,218,255 & 21,869    & 98.2\% \\
8  & 2,476,004 & 94,161    & 96.2\% \\
12 & 5,214,537 & 1,369,197 & 73.7\% \\
15 & 3,821,620 & 1,915,186 & 49.9\% \\
20 & 802,154   & 648,008   & 19.2\% \\
30 & 2,705     & 2,695     & 0.4\% \\
43 & 1         & 1         & 0\% \\
\bottomrule
\end{tabular}
\end{table}

Low vertex-count polytopes exhibit massive duplication (a typical
6-vertex polytope is generated by $\sim 114$ distinct CWS), while
high vertex-count polytopes are nearly all unique.

% ══════════════════════════════════════════════════════════════════════════
\section{Formal Verification}
\label{sec:verification}

Given the scale of the computation ($\sim 1.83 \times 10^{11}$ records)
and the impossibility of manually inspecting the output, we employ formal
verification to establish mathematical guarantees about the pipeline's
correctness.  We use two complementary tools: CBMC~\cite{CBMC} (bounded
model checking for C/C++) and Frama-C~\cite{FramaC} (abstract
interpretation for C).

\subsection{Verified properties}

Ten verification harnesses establish the following properties:

\begin{table}[h]
\centering
\caption{Formal verification harnesses and their targets.}
\label{tab:harnesses}
\small
\begin{tabular}{@{}llp{7cm}@{}}
\toprule
ID & Target & Property \\
\midrule
1.1 & \texttt{key\_less} & Strict weak ordering (irreflexivity, asymmetry,
      transitivity, equivalence consistency) \\
1.2 & \texttt{Hash128Hasher} & Consistency with \texttt{operator==} \\
1.3 & \texttt{hash\_normal\_form} & Correct byte count and array bounds \\
1.4 & \texttt{MergeRecord} & Binary layout: 72 bytes, no padding, correct
      offsets \\
1.5 & Two-way merge & Count preservation: $\sum \mathit{out} = \sum \mathit{acc}
      + \sum \mathit{shard}$ \\
1.6 & $K$-way merge & Count preservation across $K$ streams \\
1.7 & \texttt{SortedBinaryReader} & State machine: visits every record
      exactly once, in order \\
1.9 & Accounting & $\mathit{processed} = \mathit{unique} + \mathit{duplicate}
      + \mathit{failed}$ \\
2.5 & \texttt{Sort\_VL} comparator & Total order on vertex indices (no
      ties for distinct elements) \\
2.6 & \texttt{palp\_compute\_nf} & Correct CWS struct population and
      pipeline orchestration \\
\bottomrule
\end{tabular}
\end{table}

\subsection{Normal form determinism}

The most critical correctness property is that the normal form hash is
a deterministic function of the input weights.  Non-determinism would
cause the same polytope to receive different hashes on different runs,
inflating the unique polytope count.

The primary risk identified during planning was \texttt{qsort}
non-determinism: the C standard does not require \texttt{qsort} to be
stable, so if the comparator admits ties, the sorted order could vary
between calls.  We resolve this by proving (via CBMC, harness~2.5) that
the \texttt{diff} comparator used by PALP's \texttt{Sort\_VL}:
\begin{enumerate}
  \item Produces no integer overflow for vertex indices in $[0, 63]$.
  \item Is antisymmetric: $\mathrm{diff}(a,b) > 0 \iff \mathrm{diff}(b,a) < 0$.
  \item Is transitive.
  \item Admits \emph{no ties} among distinct elements:
        $\mathrm{diff}(a,b) = 0 \implies a = b$.
\end{enumerate}
Property~4 is the key result: since vertex indices within a
\texttt{VertexNumList} are distinct integers, the comparator never
returns zero for distinct elements.  Therefore \texttt{qsort}'s output
is uniquely determined, and the full chain
\[
  \text{weights} \xrightarrow{\text{Make\_CWS\_Points}} P
  \xrightarrow{\text{Find\_Equations}} (V, E)
  \xrightarrow{\text{Sort\_VL}} V'
  \xrightarrow{\text{Make\_Poly\_Sym\_NF}} \NF
  \xrightarrow{\text{xxHash128}} h
\]
is deterministic.

\subsection{Frama-C abstract interpretation of PALP}

We ran Frama-C's Eva plugin on the PALP C source code to detect potential
buffer overflows, uninitialized reads, and integer overflows.  Key
findings:
\begin{itemize}
  \item The lattice point count \texttt{P->np} is bounded within
        $[1, 2{,}000{,}000]$, confirming that the \texttt{POINT\_Nmax}
        array bound is respected.
  \item With \texttt{PALP\_FAST\_ASSERT} enabled: \textbf{zero sure
        alarms} on the critical computation path.
  \item With standard \texttt{NDEBUG}: \textbf{one sure alarm}
        (uninitialized read at \texttt{Vertex.c:598}), confirming that
        \texttt{PALP\_FAST\_ASSERT} is load-bearing for correctness.
\end{itemize}

\subsection{Runtime accounting invariant}

As a defence-in-depth measure, the pipeline checks the accounting
identity
\[
  \mathit{processed} = \mathit{unique} + \mathit{duplicate} + \mathit{failed}
\]
at the end of each run.  This identity held exactly across the full
$\sim 16.7 \times 10^9$ records in the processed subset, providing
empirical confirmation that no records were silently lost or
double-counted.

% ══════════════════════════════════════════════════════════════════════════
\section{Results}
\label{sec:results}

% TODO: Fill in with final results:
% - Total number of unique 5D reflexive polytopes found
% - Total CWS processed
% - Duplication rate (overall and by vertex count)
% - Hodge number distributions
% - Comparison with known results (lower-dimensional checks)
% - Any notable or surprising polytopes
% - Mirror symmetry verification: h^{1,1}(X) = h^{d-2,1}(X^*)

% ══════════════════════════════════════════════════════════════════════════
\section{Discussion}
\label{sec:discussion}

\subsection{Completeness and correctness guarantees}

The correctness of our classification rests on three pillars:
\begin{enumerate}
  \item \textbf{PALP's normal form algorithm} correctly computes the
        canonical representative under $\GL(5,\ZZ)$ equivalence.  This
        algorithm has been used extensively in the mathematical physics
        community since its introduction~\cite{PALP} and was validated
        against known results in dimensions 2--4.
  \item \textbf{Formal verification} of the software pipeline establishes
        that the implementation correctly orchestrates the computation:
        sorting is correct, deduplication preserves counts, checkpoint
        I/O is faithful, and the normal form hash is deterministic.
  \item \textbf{The accounting invariant} ($\mathit{processed} =
        \mathit{unique} + \mathit{duplicate} + \mathit{failed}$) provides
        a global consistency check that is extremely sensitive to any
        systematic error.
\end{enumerate}

The residual risks are:
\begin{itemize}
  \item \emph{Hash collisions:} With $\sim 5 \times 10^8$ unique polytopes
        and a 128-bit hash, the expected number of collisions is $\sim
        10^{-3}$.  This is negligible but could be eliminated entirely by
        a post-processing pass that compares full normal forms within hash
        buckets.
  \item \emph{PALP mathematical bugs:} A bug in PALP's normal form
        algorithm could cause two equivalent polytopes to receive different
        normal forms.  This would result in overcounting.  We consider this
        unlikely given PALP's extensive use and validation, but it cannot
        be ruled out by software verification alone.
\end{itemize}

\subsection{Comparison with previous work}

% TODO: Compare with:
% - Kreuzer-Skarke 4D classification (methodology, scale)
% - Any partial 5D results in the literature
% - Anderson et al., He et al. — machine learning on polytope databases

\subsection{Applications}

The resulting database enables:
\begin{itemize}
  \item Systematic study of the landscape of 5D Calabi--Yau threefolds
        from toric hypersurfaces.
  \item Computation of Hodge number distributions and mirror symmetry
        verification for the 5D case.
  \item Training data for machine learning approaches to string
        phenomenology~\cite{He2017}.
  \item Identification of extremal polytopes (maximum vertex count, maximum
        lattice points, etc.) in dimension~5.
\end{itemize}

% ══════════════════════════════════════════════════════════════════════════
\section{Conclusion}
\label{sec:conclusion}

We have presented a complete computational pipeline for classifying
five-dimensional reflexive polytopes from combined weight systems.
The pipeline processes $\sim 183$~billion CWS records ($\sim 3.2$~TB),
computing normal forms via PALP and deduplicating via xxHash128 and
external sort-merge.  Formal verification using CBMC and Frama-C
provides mathematical guarantees of correctness for all critical code
paths.

% TODO: State final polytope count and key findings.

The code and verification harnesses are publicly available at
\url{https://github.com/...}.

% ══════════════════════════════════════════════════════════════════════════
\appendix
\section{Canonical Overlap Table}
\label{app:5d-structures}

For each overlap hypergraph we display one canonical Corollary~1 ordering
of good simplices $S_1,\dots,S_\lambda$, the resulting ordered new-vertex
profile, and the overlap matrix $O_{ij} = |S_i \cap S_j|$.  If a given
hypergraph admits several Corollary~1 orderings, only the lexicographically
first one is shown here.

\begingroup
\scriptsize
\setlength{\tabcolsep}{4pt}
\renewcommand{\arraystretch}{1.08}
\begin{longtable}{@{}r>{\raggedright\arraybackslash}p{1.5cm}>{\raggedright\arraybackslash}p{8.4cm}>{\raggedright\arraybackslash}p{4.2cm}@{}}
\caption{Canonical overlap data for the $47$ combinatorial types of $5$-dimensional minimal polytopes.}
\label{tab:5d-structures}\\
\toprule
ID & Profile & Good simplices & Overlap matrix \\
\midrule
\endfirsthead
\toprule
ID & Profile & Good simplices & Overlap matrix \\
\midrule
\endhead
\bottomrule
\endfoot
1 & $[6]$ & $S_{1}=\{1,2,3,4,5,6\}$ & $\left(\begin{smallmatrix}0\end{smallmatrix}\right)$ \\
2 & $[5,2]$ & $S_{1}=\{1,2\}, S_{2}=\{3,4,5,6,7\}$ & $\left(\begin{smallmatrix}0 & 0\\0 & 0\end{smallmatrix}\right)$ \\
3 & $[5,2]$ & $S_{1}=\{1,2,3,4,5\}, S_{2}=\{1,2,3,6,7\}$ & $\left(\begin{smallmatrix}0 & 3\\3 & 0\end{smallmatrix}\right)$ \\
4 & $[4,3]$ & $S_{1}=\{1,2,3\}, S_{2}=\{1,4,5,6,7\}$ & $\left(\begin{smallmatrix}0 & 1\\1 & 0\end{smallmatrix}\right)$ \\
5 & $[4,3]$ & $S_{1}=\{1,2,3\}, S_{2}=\{4,5,6,7\}$ & $\left(\begin{smallmatrix}0 & 0\\0 & 0\end{smallmatrix}\right)$ \\
6 & $[4,3]$ & $S_{1}=\{1,2,3,4\}, S_{2}=\{1,2,5,6,7\}$ & $\left(\begin{smallmatrix}0 & 2\\2 & 0\end{smallmatrix}\right)$ \\
7 & $[4,3]$ & $S_{1}=\{1,2,3,4\}, S_{2}=\{1,5,6,7\}$ & $\left(\begin{smallmatrix}0 & 1\\1 & 0\end{smallmatrix}\right)$ \\
8 & $[4,2,2]$ & $S_{1}=\{1,2\}, S_{2}=\{1,2,3,4\}, S_{3}=\{5,6,7,8\}$ & $\left(\begin{smallmatrix}0 & 2 & 0\\2 & 0 & 0\\0 & 0 & 0\end{smallmatrix}\right)$ \\
9 & $[4,2,2]$ & $S_{1}=\{1,2\}, S_{2}=\{3,4\}, S_{3}=\{5,6,7,8\}$ & $\left(\begin{smallmatrix}0 & 0 & 0\\0 & 0 & 0\\0 & 0 & 0\end{smallmatrix}\right)$ \\
10 & $[4,2,2]$ & $S_{1}=\{1,2\}, S_{2}=\{3,4,5,6\}, S_{3}=\{3,4,7,8\}$ & $\left(\begin{smallmatrix}0 & 0 & 0\\0 & 0 & 2\\0 & 2 & 0\end{smallmatrix}\right)$ \\
11 & $[4,2,2]$ & $S_{1}=\{1,2,3,4\}, S_{2}=\{1,2,5,6\}, S_{3}=\{1,2,7,8\}$ & $\left(\begin{smallmatrix}0 & 2 & 2\\2 & 0 & 2\\2 & 2 & 0\end{smallmatrix}\right)$ \\
12 & $[4,2,2]$ & $S_{1}=\{1,2,3,4\}, S_{2}=\{1,2,5,6\}, S_{3}=\{1,3,4,7,8\}$ & $\left(\begin{smallmatrix}0 & 2 & 3\\2 & 0 & 1\\3 & 1 & 0\end{smallmatrix}\right)$ \\
13 & $[4,2,2]$ & $S_{1}=\{1,2,3,4\}, S_{2}=\{1,2,5,6\}, S_{3}=\{2,3,4,7,8\}$ & $\left(\begin{smallmatrix}0 & 2 & 3\\2 & 0 & 1\\3 & 1 & 0\end{smallmatrix}\right)$ \\
14 & $[4,2,2]$ & $S_{1}=\{1,2,3,4\}, S_{2}=\{1,2,5,6\}, S_{3}=\{3,4,7,8\}$ & $\left(\begin{smallmatrix}0 & 2 & 2\\2 & 0 & 0\\2 & 0 & 0\end{smallmatrix}\right)$ \\
15 & $[3,3,2]$ & $S_{1}=\{1,2\}, S_{2}=\{1,2,3,4,5\}, S_{3}=\{3,6,7,8\}$ & $\left(\begin{smallmatrix}0 & 2 & 0\\2 & 0 & 1\\0 & 1 & 0\end{smallmatrix}\right)$ \\
16 & $[3,3,2]$ & $S_{1}=\{1,2\}, S_{2}=\{3,4,5\}, S_{3}=\{3,6,7,8\}$ & $\left(\begin{smallmatrix}0 & 0 & 0\\0 & 0 & 1\\0 & 1 & 0\end{smallmatrix}\right)$ \\
17 & $[3,3,2]$ & $S_{1}=\{1,2\}, S_{2}=\{3,4,5\}, S_{3}=\{4,6,7,8\}$ & $\left(\begin{smallmatrix}0 & 0 & 0\\0 & 0 & 1\\0 & 1 & 0\end{smallmatrix}\right)$ \\
18 & $[3,3,2]$ & $S_{1}=\{1,2\}, S_{2}=\{3,4,5\}, S_{3}=\{5,6,7,8\}$ & $\left(\begin{smallmatrix}0 & 0 & 0\\0 & 0 & 1\\0 & 1 & 0\end{smallmatrix}\right)$ \\
19 & $[3,3,2]$ & $S_{1}=\{1,2\}, S_{2}=\{3,4,5\}, S_{3}=\{6,7,8\}$ & $\left(\begin{smallmatrix}0 & 0 & 0\\0 & 0 & 0\\0 & 0 & 0\end{smallmatrix}\right)$ \\
20 & $[3,3,2]$ & $S_{1}=\{1,2,3\}, S_{2}=\{1,2,3,4,5\}, S_{3}=\{1,6,7,8\}$ & $\left(\begin{smallmatrix}0 & 3 & 1\\3 & 0 & 1\\1 & 1 & 0\end{smallmatrix}\right)$ \\
21 & $[3,3,2]$ & $S_{1}=\{1,2,3\}, S_{2}=\{1,2,3,4,5\}, S_{3}=\{6,7,8\}$ & $\left(\begin{smallmatrix}0 & 3 & 0\\3 & 0 & 0\\0 & 0 & 0\end{smallmatrix}\right)$ \\
22 & $[3,3,2]$ & $S_{1}=\{1,2,3\}, S_{2}=\{1,4,5\}, S_{3}=\{1,6,7,8\}$ & $\left(\begin{smallmatrix}0 & 1 & 1\\1 & 0 & 1\\1 & 1 & 0\end{smallmatrix}\right)$ \\
23 & $[3,3,2]$ & $S_{1}=\{1,2,3\}, S_{2}=\{1,4,5\}, S_{3}=\{6,7,8\}$ & $\left(\begin{smallmatrix}0 & 1 & 0\\1 & 0 & 0\\0 & 0 & 0\end{smallmatrix}\right)$ \\
24 & $[3,3,2]$ & $S_{1}=\{1,2,3\}, S_{2}=\{1,4,5,6\}, S_{3}=\{1,4,7,8\}$ & $\left(\begin{smallmatrix}0 & 1 & 1\\1 & 0 & 2\\1 & 2 & 0\end{smallmatrix}\right)$ \\
25 & $[3,3,2]$ & $S_{1}=\{1,2,3\}, S_{2}=\{1,4,5,6\}, S_{3}=\{2,3,4,7,8\}$ & $\left(\begin{smallmatrix}0 & 1 & 2\\1 & 0 & 1\\2 & 1 & 0\end{smallmatrix}\right)$ \\
26 & $[3,3,2]$ & $S_{1}=\{1,2,3\}, S_{2}=\{1,4,5,6\}, S_{3}=\{2,3,5,7,8\}$ & $\left(\begin{smallmatrix}0 & 1 & 2\\1 & 0 & 1\\2 & 1 & 0\end{smallmatrix}\right)$ \\
27 & $[3,3,2]$ & $S_{1}=\{1,2,3\}, S_{2}=\{1,4,5,6\}, S_{3}=\{2,3,6,7,8\}$ & $\left(\begin{smallmatrix}0 & 1 & 2\\1 & 0 & 1\\2 & 1 & 0\end{smallmatrix}\right)$ \\
28 & $[3,3,2]$ & $S_{1}=\{1,2,3\}, S_{2}=\{1,4,5,6\}, S_{3}=\{2,3,7,8\}$ & $\left(\begin{smallmatrix}0 & 1 & 2\\1 & 0 & 0\\2 & 0 & 0\end{smallmatrix}\right)$ \\
29 & $[3,3,2]$ & $S_{1}=\{1,2,3\}, S_{2}=\{1,4,5,6\}, S_{3}=\{4,5,6,7,8\}$ & $\left(\begin{smallmatrix}0 & 1 & 0\\1 & 0 & 3\\0 & 3 & 0\end{smallmatrix}\right)$ \\
30 & $[3,3,2]$ & $S_{1}=\{1,2,3\}, S_{2}=\{1,4,5,6\}, S_{3}=\{4,7,8\}$ & $\left(\begin{smallmatrix}0 & 1 & 0\\1 & 0 & 1\\0 & 1 & 0\end{smallmatrix}\right)$ \\
31 & $[3,3,2]$ & $S_{1}=\{1,2,3\}, S_{2}=\{2,4,5,6\}, S_{3}=\{4,7,8\}$ & $\left(\begin{smallmatrix}0 & 1 & 0\\1 & 0 & 1\\0 & 1 & 0\end{smallmatrix}\right)$ \\
32 & $[3,3,2]$ & $S_{1}=\{1,2,3\}, S_{2}=\{3,4,5,6\}, S_{3}=\{4,7,8\}$ & $\left(\begin{smallmatrix}0 & 1 & 0\\1 & 0 & 1\\0 & 1 & 0\end{smallmatrix}\right)$ \\
33 & $[3,2,2,2]$ & $S_{1}=\{1,2\}, S_{2}=\{1,2,3,4\}, S_{3}=\{5,6\}, S_{4}=\{7,8,9\}$ & $\left(\begin{smallmatrix}0 & 2 & 0 & 0\\2 & 0 & 0 & 0\\0 & 0 & 0 & 0\\0 & 0 & 0 & 0\end{smallmatrix}\right)$ \\
34 & $[3,2,2,2]$ & $S_{1}=\{1,2\}, S_{2}=\{1,2,3,4\}, S_{3}=\{5,6,7\}, S_{4}=\{5,8,9\}$ & $\left(\begin{smallmatrix}0 & 2 & 0 & 0\\2 & 0 & 0 & 0\\0 & 0 & 0 & 1\\0 & 0 & 1 & 0\end{smallmatrix}\right)$ \\
35 & $[3,2,2,2]$ & $S_{1}=\{1,2\}, S_{2}=\{1,2,3,4,5\}, S_{3}=\{3,6,7\}, S_{4}=\{3,8,9\}$ & $\left(\begin{smallmatrix}0 & 2 & 0 & 0\\2 & 0 & 1 & 1\\0 & 1 & 0 & 1\\0 & 1 & 1 & 0\end{smallmatrix}\right)$ \\
36 & $[3,2,2,2]$ & $S_{1}=\{1,2\}, S_{2}=\{1,2,3,4,5\}, S_{3}=\{3,6,7\}, S_{4}=\{8,9\}$ & $\left(\begin{smallmatrix}0 & 2 & 0 & 0\\2 & 0 & 1 & 0\\0 & 1 & 0 & 0\\0 & 0 & 0 & 0\end{smallmatrix}\right)$ \\
37 & $[3,2,2,2]$ & $S_{1}=\{1,2\}, S_{2}=\{3,4\}, S_{3}=\{5,6\}, S_{4}=\{7,8,9\}$ & $\left(\begin{smallmatrix}0 & 0 & 0 & 0\\0 & 0 & 0 & 0\\0 & 0 & 0 & 0\\0 & 0 & 0 & 0\end{smallmatrix}\right)$ \\
38 & $[3,2,2,2]$ & $S_{1}=\{1,2\}, S_{2}=\{3,4\}, S_{3}=\{5,6,7\}, S_{4}=\{5,6,7,8,9\}$ & $\left(\begin{smallmatrix}0 & 0 & 0 & 0\\0 & 0 & 0 & 0\\0 & 0 & 0 & 3\\0 & 0 & 3 & 0\end{smallmatrix}\right)$ \\
39 & $[3,2,2,2]$ & $S_{1}=\{1,2\}, S_{2}=\{3,4\}, S_{3}=\{5,6,7\}, S_{4}=\{5,8,9\}$ & $\left(\begin{smallmatrix}0 & 0 & 0 & 0\\0 & 0 & 0 & 0\\0 & 0 & 0 & 1\\0 & 0 & 1 & 0\end{smallmatrix}\right)$ \\
40 & $[3,2,2,2]$ & $S_{1}=\{1,2\}, S_{2}=\{3,4,5\}, S_{3}=\{3,4,5,6,7\}, S_{4}=\{3,8,9\}$ & $\left(\begin{smallmatrix}0 & 0 & 0 & 0\\0 & 0 & 3 & 1\\0 & 3 & 0 & 1\\0 & 1 & 1 & 0\end{smallmatrix}\right)$ \\
41 & $[3,2,2,2]$ & $S_{1}=\{1,2\}, S_{2}=\{3,4,5\}, S_{3}=\{3,6,7\}, S_{4}=\{3,8,9\}$ & $\left(\begin{smallmatrix}0 & 0 & 0 & 0\\0 & 0 & 1 & 1\\0 & 1 & 0 & 1\\0 & 1 & 1 & 0\end{smallmatrix}\right)$ \\
42 & $[3,2,2,2]$ & $S_{1}=\{1,2\}, S_{2}=\{3,4,5\}, S_{3}=\{3,6,7\}, S_{4}=\{4,5,8,9\}$ & $\left(\begin{smallmatrix}0 & 0 & 0 & 0\\0 & 0 & 1 & 2\\0 & 1 & 0 & 0\\0 & 2 & 0 & 0\end{smallmatrix}\right)$ \\
43 & $[3,2,2,2]$ & $S_{1}=\{1,2,3\}, S_{2}=\{1,2,3,4,5\}, S_{3}=\{1,6,7\}, S_{4}=\{1,8,9\}$ & $\left(\begin{smallmatrix}0 & 3 & 1 & 1\\3 & 0 & 1 & 1\\1 & 1 & 0 & 1\\1 & 1 & 1 & 0\end{smallmatrix}\right)$ \\
44 & $[3,2,2,2]$ & $S_{1}=\{1,2,3\}, S_{2}=\{1,4,5\}, S_{3}=\{1,6,7\}, S_{4}=\{1,8,9\}$ & $\left(\begin{smallmatrix}0 & 1 & 1 & 1\\1 & 0 & 1 & 1\\1 & 1 & 0 & 1\\1 & 1 & 1 & 0\end{smallmatrix}\right)$ \\
45 & $[3,2,2,2]$ & $S_{1}=\{1,2,3\}, S_{2}=\{1,4,5\}, S_{3}=\{1,6,7\}, S_{4}=\{2,3,8,9\}$ & $\left(\begin{smallmatrix}0 & 1 & 1 & 2\\1 & 0 & 1 & 0\\1 & 1 & 0 & 0\\2 & 0 & 0 & 0\end{smallmatrix}\right)$ \\
46 & $[2,2,2,2,2]$ & $S_{1}=\{1,2\}, S_{2}=\{1,2,3,4\}, S_{3}=\{5,6\}, S_{4}=\{7,8\}, S_{5}=\{9,10\}$ & $\left(\begin{smallmatrix}0 & 2 & 0 & 0 & 0\\2 & 0 & 0 & 0 & 0\\0 & 0 & 0 & 0 & 0\\0 & 0 & 0 & 0 & 0\\0 & 0 & 0 & 0 & 0\end{smallmatrix}\right)$ \\
47 & $[2,2,2,2,2]$ & $S_{1}=\{1,2\}, S_{2}=\{3,4\}, S_{3}=\{5,6\}, S_{4}=\{7,8\}, S_{5}=\{9,10\}$ & $\left(\begin{smallmatrix}0 & 0 & 0 & 0 & 0\\0 & 0 & 0 & 0 & 0\\0 & 0 & 0 & 0 & 0\\0 & 0 & 0 & 0 & 0\\0 & 0 & 0 & 0 & 0\end{smallmatrix}\right)$ \\
\end{longtable}
\endgroup

% ══════════════════════════════════════════════════════════════════════════
\section*{Acknowledgements}

% TODO

% ══════════════════════════════════════════════════════════════════════════
\bibliographystyle{unsrtnat}
% Inline bib for draft — replace with .bib file for submission
\begin{thebibliography}{99}

\bibitem{Batyrev1994}
V.~V. Batyrev,
\emph{Dual polyhedra and mirror symmetry for Calabi--Yau hypersurfaces
in toric varieties},
J.~Algebraic Geom. \textbf{3} (1994) 493--535.
\href{https://arxiv.org/abs/alg-geom/9310003}{arXiv:alg-geom/9310003}.

\bibitem{BatyrevBorisov1996}
V.~V. Batyrev and L.~A. Borisov,
\emph{On Calabi--Yau complete intersections in toric varieties},
in: Higher-Dimensional Complex Varieties (Trento, 1994), de~Gruyter,
Berlin, 1996, pp.~39--65.
\href{https://arxiv.org/abs/alg-geom/9412017}{arXiv:alg-geom/9412017}.

\bibitem{Koelman1991}
R.~J. Koelman,
\emph{The number of moduli of families of curves on toric surfaces},
Ph.D. thesis, Katholieke Universiteit Nijmegen, 1991.

\bibitem{KreuzerSkarke1998}
M.~Kreuzer and H.~Skarke,
\emph{Classification of reflexive polyhedra in three dimensions},
Adv.~Theor.~Math.~Phys. \textbf{2} (1998) 853--871.
\href{https://arxiv.org/abs/hep-th/9805190}{arXiv:hep-th/9805190}.

\bibitem{KreuzerSkarke2000}
M.~Kreuzer and H.~Skarke,
\emph{Complete classification of reflexive polyhedra in four dimensions},
Adv.~Theor.~Math.~Phys. \textbf{4} (2000) 1209--1230.
\href{https://arxiv.org/abs/hep-th/0002240}{arXiv:hep-th/0002240}.

\bibitem{KreuzerSkarke1997}
M.~Kreuzer and H.~Skarke,
\emph{On the classification of reflexive polyhedra},
Commun.~Math.~Phys. \textbf{185} (1997) 495--508.
\href{https://arxiv.org/abs/hep-th/9512204}{arXiv:hep-th/9512204}.

\bibitem{KreuzerSkarkeCYcsws}
M.~Kreuzer and H.~Skarke,
\emph{Classification scheme and weight systems},
online classification page, accessed 2026.
\url{https://hep.itp.tuwien.ac.at/~kreuzer/CY/CYcsws.html}.

\bibitem{PALP}
M.~Kreuzer and H.~Skarke,
\emph{PALP: A Package for Analysing Lattice Polytopes with applications
to toric geometry},
Comput.~Phys.~Commun. \textbf{157} (2004) 87--106.
\href{https://arxiv.org/abs/math/0204356}{arXiv:math/0204356}.

\bibitem{xxHash}
Y.~Collet,
\emph{xxHash — Extremely fast non-cryptographic hash algorithm},
\url{https://github.com/Cyan4973/xxHash}, 2012--2024.

\bibitem{CBMC}
D.~Kroening and M.~Tautschnig,
\emph{CBMC -- C Bounded Model Checker},
in: TACAS 2014, LNCS 8413, Springer, 2014, pp.~389--391.

\bibitem{FramaC}
P.~Cuoq, F.~Kirchner, N.~Kosmatov, V.~Prevosto, J.~Signoles, and
B.~Yakobowski,
\emph{Frama-C: A software analysis perspective},
in: SEFM 2012, LNCS 7504, Springer, 2012, pp.~233--247.

\bibitem{He2017}
Y.-H. He,
\emph{Machine-learning the string landscape},
Phys.~Lett.~B \textbf{774} (2017) 564--568.
\href{https://arxiv.org/abs/1706.02714}{arXiv:1706.02714}.

\end{thebibliography}

\end{document}
